%======================================================================
\section{Requisitos}\label{sec:requisitos}
%======================================================================
Os requisitos abaixo revisam os do Projeto 1 do primeiro semestre, atualizados com as
decisões tomadas desde então (Seção~\ref{sec:evolucao}). O escopo desta fase é o quarto de
um residente: portas, presença, temperatura e a pulseira. Outros cômodos e as áreas comuns
ficam para depois.

\subsection{Funcionais}

\begin{table}[H]
\centering
\caption{Requisitos funcionais.}
\label{tab:rf}
\small
\begin{tabular}{@{}l>{\raggedright\arraybackslash}p{12.2cm}@{}}
\toprule
ID & Requisito \\
\midrule
\multicolumn{2}{@{}l}{\textit{Pulseira}} \\
RF01 & Adquirir aceleração e velocidade angular nos 3 eixos a 200\,Hz, com faixa de $\pm$8\,g \\
RF02 & Detectar a queda na própria pulseira: gatilho em picos de $\amag \ge \gatilhoLon$\,g e classificador sobre \nAtributos{} atributos da janela (modelo provisório até retreinar com dados do pulso) \\
RF03 & Enviar a queda ao hub imediatamente, reenviando até a confirmação \\
RF04 & Depois de uma queda, não gerar outra por 30\,s \\
RF05 & Cancelar a queda ou o pedido de ajuda com dois toques rápidos no botão (``estou bem''), nos 15\,s seguintes; um aperto simples não faz nada \\
RF20 & Pedir ajuda segurando o botão por 3\,s, fora de alerta \\
RF06 & Enviar um sinal de vida a cada 30\,s, com a bateria e o nível de atividade \\
RF07 & Ter um modo de coleta que grava as amostras brutas, para formar o conjunto de dados \\
\midrule
\multicolumn{2}{@{}l}{\textit{Hub do quarto}} \\
RF08 & Detectar toda transição de cada porta (quantidade configurável) e publicar o estado \\
RF09 & Detectar presença e publicar o estado \\
RF10 & Medir a temperatura a cada 60\,s (16 amostras, curva de calibração da cadeia completa) e publicá-la quando mudar 0,2\,°C ou a cada 5\,min \\
RF11 & Executar as regras de tempo e publicar os episódios de alerta \\
RF12 & Sinalizar o estado e os alertas por um LED \\
RF13 & Aplicar parâmetros recebidos por MQTT sem recompilar, guardando-os na memória não volátil \\
RF14 & Publicar no barramento CAN os estados dos sensores locais \\
RF15 & Sincronizar a hora por SNTP \\
\midrule
\multicolumn{2}{@{}l}{\textit{Central}} \\
RF16 & Broker MQTT com autenticação, disponível continuamente na rede local \\
RF17 & Gravar todos os tópicos num banco de séries temporais, com retenção de 30 dias ou mais \\
RF18 & Painel com o estado de todos os quartos, o histórico das portas, a temperatura, a atividade e os alertas \\
RF19 & Notificar cada alerta ao cuidador pelo Telegram: tipo, quarto e horário \\
\bottomrule
\end{tabular}
\end{table}

\subsection{Não funcionais}

\begin{table}[H]
\centering
\caption{Requisitos não funcionais, com as metas e como medi-las.}
\label{tab:rnf}
\small
\begin{tabular}{@{}l>{\raggedright\arraybackslash}p{3.6cm}>{\raggedright\arraybackslash}p{4.4cm}>{\raggedright\arraybackslash}p{4.0cm}@{}}
\toprule
ID & Requisito & Meta & Como medir \\
\midrule
RNF01 & Amostragem estável & 200\,Hz, \textit{jitter} $\le \pm$5\,ms & intervalos no modo de coleta \\
RNF02 & Latência da queda & pico $\to$ MQTT $\le$ 500\,ms & hora na pulseira, no hub e no broker \\
RNF03 & Latência da porta & transição $\to$ MQTT $\le$ 200\,ms & log com hora nos dois lados \\
RNF04 & Latência da notificação & alerta $\to$ Telegram $\le$ 5\,s & hora do alerta e da mensagem \\
RNF05 & Disponibilidade & 24/7, reconexão automática com espera crescente & 24\,h derrubando a rede \\
RNF06 & Sem perda de eventos & nenhum evento QoS 1 perdido & derrubar o Wi-Fi durante eventos \\
RNF07 & Watchdog & reinício se uma task crítica travar por mais de 30\,s & travamento forçado \\
RNF08 & Detecção de queda & sensibilidade $\ge$ 90\,\%; $\le$ 1 alarme falso por dia por residente & ensaios e uso contínuo \\
RNF09 & Segurança & MQTT com usuário, senha e ACL; ESP-NOW cifrado; segredos fora do repositório & revisão da configuração \\
RNF10 & Temperatura & resolução $\le$ 0,1\,°C; exatidão $\pm$1\,°C & comparação com termopar \\
RNF11 & Sem internet & tudo funciona na rede local, exceto o Telegram & ensaio com a internet desligada \\
RNF13 & Diagnóstico & log na serial: tasks, conexões, eventos & inspeção \\
RNF14 & Autonomia da pulseira & a definir depois de medir o consumo & medição com bateria \\
\bottomrule
\end{tabular}
\end{table}

\subsection{Restrições}

\begin{itemize}
  \item \textbf{Processamento na ESP32} (diretriz da disciplina): a detecção de queda roda na
        pulseira; as regras e os alertas, no hub. A central só armazena, mostra e notifica.
  \item Firmware em C com ESP-IDF 6 e FreeRTOS, na ESP32-S3-DevKitC-1.
  \item Hub alimentado em 5\,V; pulseira a bateria (no protótipo, um \textit{power bank}).
  \item Central com ferramentas abertas: Mosquitto, Telegraf, InfluxDB, Grafana e Python.
  \item Integração com o barramento CAN da sala, compartilhado com os outros grupos.
\end{itemize}
